% Part: normal-modal-logic
% Chapter: completeness
% Section: modalities-ccs

\documentclass[../../../include/open-logic-section]{subfiles}

\begin{document}

\olfileid{nml}{com}{mod}

\olsection{模态词与完全一致集}

\begin{explain}
当我们构造一个模型~$\mModel{M^\Sigma}$，其可能世界集由某个正规模态逻辑~$\Sigma$ 的完全 $\Sigma$ 一致集~$\Delta$ 给出时，我们还需要在这些“世界”之间定义一个可达关系~$R^\Sigma$。我们希望将可达关系（以及赋值 $V^\Sigma$）定义成满足：$\mSat{M^\Sigma}{!A}[\Delta]$ 当且仅当 $!A \in \Delta$。我们应当怎么做呢？

一旦定义了可达关系，“在世界中为真”的定义就保证了 \iftag{prvBox}{$\mSat{M^\Sigma}{\Box!A}[\Delta]$当且仅当对所有满足 $R^\Sigma\Delta\Delta'$ 的 $\Delta'$ 都有$\mSat{M^\Sigma}{!A}[\Delta']$}{$\mSat{M^\Sigma}{\Diamond!A}[\Delta]$当且仅当至少存在一个满足 $R^\Sigma\Delta\Delta'$ 的 $\Delta'$ 使得$\mSat{M^\Sigma}{!A}[\Delta']$}。要证明$\mSat{M^\Sigma}{!A}[\Delta]$当且仅当 $!A \in \Delta$，特别地需要该结论对以模态算子开头的公式成立，即 \iftag{prvBox}{$\mSat{M^\Sigma}{\Box!A}[\Delta]$当且仅当 $\Box !A \in \Delta$}{$\mSat{M^\Sigma}{\Diamond!A}[\Delta]$当且仅当 $\Diamond !A \in \Delta$}。将此要求与 \iftag{prvBox}{$\Box !A$}{$\Diamond !A$} 在世界中为真的定义相结合，得到：
\[
\iftag{prvBox}{%
  \Box!A \in \Delta \text{ 当且仅当 } !A \in \Delta' \text{ 对所有满足 } R^\Sigma\Delta\Delta'}{%
  \Diamond!A \in \Delta \text{ 的 $\Delta'$ 成立 } !A \in \Delta' \text{ 当且仅当对至少一个满足 } R^\Sigma\Delta\Delta'}
\]
的 $\Delta'$ 成立
\iftag{prvBox}{考虑从左到右的方向：它说的是，若 $\Box !A \in \Delta$，则对任意 $!A$ 及任意满足 $R^\Sigma\Delta\Delta'$ 的 $\Delta'$，都有 $!A \in \Delta'$。如果我们规定 $R^\Sigma\Delta\Delta'$ 当且仅当：对所有的 $\Box!A \in \Delta$，都有 $!A \in \Delta'$，那么这就成立了。我们可以将“当且仅当”右边的条件更简洁地写为：$\Setabs{!A}{\Box!A \in \Delta}
  \subseteq \Delta'$。}{考虑从右到左的方向：它说的是，对任意 $!A$ 及满足 $R^\Sigma\Delta\Delta'$ 的 $\Delta'$，若 $!A \in \Delta'$ 则 $\Diamond!A \in \Delta$。如果我们规定，只要对所有的 $!A \in \Delta'$ 都有 $\Diamond!A \in \Delta$，就成立 $R^\Sigma\Delta\Delta'$，那么这就成立了。我们可以将“当且仅当”右边的条件更简洁地写为：$\Setabs{\Diamond!A}{!A \in \Delta'} \subseteq \Delta$。}

那么问题在于：$R^\Sigma$ 的这个定义是否确实能保证 \iftag{prvBox}{$\Box !A \in \Delta$ 当且仅当$\mSat{M^\Sigma}{\Box !A}[\Delta]$？ \iftag{prvDiamond}{它是否也保证了 }{}}{}\iftag{prvDiamond}{$\Diamond !A \in \Delta$ 当且仅当$\mSat{M^\Sigma}{\Diamond !A}[\Delta]$？}{} 接下来的几个结论将确立这一点。
\end{explain}

\begin{defn}
  如果 $\Gamma$ 是一个公式集，令
  \begin{align*}
    \Box\Gamma & = \Setabs{\Box !B}{!B \in \Gamma}\\
    \Diamond\Gamma & = \Setabs{\Diamond !B}{!B \in \Gamma}\\
    \intertext{且}
    \Box^{-1}\Gamma & = \Setabs{!B}{\Box !B \in \Gamma}\\
    \Diamond^{-1}\Gamma & = \Setabs{!B}{\Diamond !B \in \Gamma}\\
  \end{align*}
\end{defn}

换句话说，$\Box\Gamma$ 是在 $\Gamma$ 的每个公式前加上 $\Box$ 所得的集合；$\Box^{-1}\Gamma$ 则是从 $\Gamma$ 中所有带 $\Box$ 的公式中去掉最前面的 $\Box$ 所得的集合。这个定义本身不太重要，但能极大地简化记号。

注意 $\Box\Box^{-1}\Gamma \subseteq \Gamma$：
\[
\Box\Box^{-1}\Gamma = \Setabs{\Box!B}{\Box!B \in \Gamma}
\]
也就是说，$\Box\Box^{-1}\Gamma$ 正好就是 $\Gamma$ 中所有以 $\Box$ 开头的公式组成的集合。

\begin{lem}\ollabel{lem:box1}
  若 $\Gamma \Proves[\Sigma] !A$，则 $\Box\Gamma \Proves[\Sigma] \Box
  !A$.
\end{lem}

\begin{proof}
  若 $\Gamma \Proves[\Sigma] !A$，则存在 $!B_1$, \dots, $!B_k
  \in \Gamma$ 使得 $\Sigma \Proves !B_1 \lif (!B_2 \lif \cdots
  (!B_n \lif !A)\cdots)$。由于 $\Sigma$ 是正规的，根据规则 \RK{},
  $\Sigma \Proves \Box!B_1 \lif (\Box!B_2 \lif \cdots (\Box!B_n \lif
  \Box!A)\cdots)$，其中显然有 $\Box!B_1$, \dots, $\Box!B_k \in
  \Box\Gamma$。因此，由定义知 $\Box\Gamma \Proves[\Sigma] \Box
  !A$.
\end{proof}

\begin{lem}\ollabel{lem:box2}
  若 $\Box^{-1} \Gamma \Proves[\Sigma] !A$，则 $\Gamma
  \Proves[\Sigma] \Box!A$.
\end{lem}

\begin{proof}
  假设 $\Box^{-1}\Gamma \Proves[\Sigma] !A$；则由
  \olref{lem:box1}，$\Box\Box^{-1}\Gamma \Proves \Box!A$。又因为
  $\Box\Box^{-1}\Gamma \subseteq \Gamma$，根据单调性，也有 $\Gamma \Proves[\Sigma]
  \Box!A$。
\end{proof}

\iftag{prvBox}{%
% Box is primitive


\begin{prop}\ollabel{prop:box}
  若 $\Gamma$ 是完全 $\Sigma$ 一致的，则 $\Box !A \in
  \Gamma$ 当且仅当：对于每个满足 $\Box^{-1} \Gamma \subseteq \Delta$ 的完全 $\Sigma$ 一致集 $\Delta$，都有 $!A \in \Delta$。
\end{prop}

\begin{proof}
  假设 $\Gamma$ 是完全 $\Sigma$ 一致的。“仅当”方向是容易的：假设 $\Box !A \in \Gamma$ 且 $\Box^{-1}\Gamma \subseteq \Delta$。由于 $\Box!A \in \Gamma$，$!A
  \in \Box^{-1}\Gamma \subseteq \Delta$，故 $!A \in \Delta$。

  对于“当”方向，我们用逆否法论证：假设
  $\Box!A \notin \Gamma$。由于 $\Gamma$ 是完全 $\Sigma$ 一致的，它是演绎封闭的，因此 $\Gamma
  \Proves/[\Sigma] \Box !A$。由 \olref{lem:box2},
  $\Box^{-1}\Gamma \Proves/[\Sigma] !A$。根据
  \olref[prf][con]{prop:consistencyfacts}\olref[prf][con]{prop:consistencyfacts-b},
  $\Box^{-1}\Gamma \cup \{ \lnot!A \}$ 是 $\Sigma$ 一致的。由 Lindenbaum 引理，存在一个完全 $\Sigma$ 一致集~$\Delta$ 使得 $\Box^{-1}\Gamma \cup \{ \lnot!A \} \subseteq
  \Delta$。由一致性，$!A \notin \Delta$。
\end{proof}


}{%
% Box is not primitive, only Diamond is

\begin{prop}\ollabel{prop:diamond}
  若 $\Gamma$ 是完全 $\Sigma$ 一致的，则 $\Diamond !A \in
  \Gamma$ 当且仅当存在某个完全 $\Sigma$ 一致集
  $\Delta$ 使得 $\Diamond\Delta \subseteq
  \Gamma$ 且 $!A \in \Delta$。
\end{prop}

\begin{proof}
  假设 $\Gamma$ 是完全 $\Sigma$ 一致的。
  从右到左的部分是简单的：设 $\Delta$ 是完全 $\Sigma$ 一致的，且满足 $\Diamond\Delta \subseteq \Gamma$ 以及 $!A \in \Delta$。
  那么 $\Diamond!A \in \Diamond\Delta \subseteq \Gamma$，即 $\Diamond!A
  \in \Gamma$。

  对于从左到右的方向，假设 $\Diamond !A \in \Gamma$。
  我们需要证明存在一个完全 $\Sigma$ 一致集 $\Delta$，满足 $!A \in \Delta$ 且 $\Diamond\Delta \subseteq \Gamma$。
  既然 $\Diamond !A \in \Gamma$，且 $\Gamma$ 是 $\Sigma$ 一致的，则 $\Box\lnot !A \notin \Gamma$。
  由于 $\Gamma$ 是完全 $\Sigma$ 一致的，它是演绎封闭的，所以 $\Gamma
  \Proves/[\Sigma] \Box\lnot !A$。
  根据 \olref{lem:box2}，$\Box^{-1}\Gamma \Proves/[\Sigma] \lnot !A$。
  由
  \olref[prf][con]{prop:consistencyfacts}\olref[prf][con]{prop:consistencyfacts-b}，
  $\Box^{-1}\Gamma \cup \{ !A \}$ 是 $\Sigma$ 一致的。
  根据 Lindenbaum 引理，存在一个完全 $\Sigma$ 一致集
  $\Delta$ 使得 $\Box^{-1}\Gamma \cup \{ !A \} \subseteq
  \Delta$。
  显然，$!A \in \Delta$。
  为证 $\Diamond\Delta \subseteq \Gamma$，假设这不成立：即存在 $!B \in \Delta$，
  使得 $\Diamond !B \notin \Gamma$。
  那么，由于 $\Gamma$ 是完全 $\Sigma$ 一致的，有 $\Box\lnot !B \in \Gamma$。
  但这样 $\lnot
  !B \in \Box^{-1}\Gamma \subseteq \Delta$ 也会成立，这将导致 $\Delta$ 不一致。
\end{proof}
}

\begin{lem}\ollabel{lem:box-iff-diamond}
  假设 $\Gamma$ 和 $\Delta$ 都是完全
  $\Sigma$ 一致的。则 $\Box^{-1}\Gamma \subseteq \Delta$ 当且仅当
  $\Diamond\Delta \subseteq \Gamma$。
\end{lem}

\begin{proof}
  “仅当”方向：假设 $\Box^{-1}\Gamma \subseteq \Delta$ 并且
  设 $\Diamond!A \in \Diamond\Delta$（即 $!A \in \Delta$）。
  为证明 $\Diamond!A \in \Gamma$，只需证明
  $\Box\lnot!A \notin \Gamma$，因为这样由极大性便有 $\lnot\Box\lnot
  !A \in \Gamma$。
  现在，若 $\Box\lnot!A \in \Gamma$，则根据假设
  $\lnot!A \in \Delta$，这与 $\Delta$ 的一致性矛盾（因为 $!A
  \in \Delta$）。
  因此 $\Box\lnot!A \notin \Gamma$，符合要求。

  “当”方向：假设 $\Diamond\Delta
  \subseteq \Gamma$。
  我们用逆否法论证：假设 $!A
  \notin \Delta$，以证明 $\Box!A \notin \Gamma$。
  如果 $!A \notin \Delta$，则由极大性有 $\lnot!A \in \Delta$，
  从而根据假设 $\Diamond\lnot!A \in \Gamma$。
  但在正规模态逻辑中，$\Diamond\lnot!A$ 等价于
  $\lnot\Box !A$，而若后者在 $\Gamma$ 中，由一致性可知
  $\Box!A \notin\Gamma$，符合要求。
\end{proof}

\iftag{prvBox}{%
  \iftag{prvDiamond}{%
% Box and Diamond are both primitive; prove
% prop:diamond using lem:box-iff-diamond


\begin{prop}\ollabel{prop:diamond}
  若 $\Gamma$ 是完全 $\Sigma$ 一致的，则 $\Diamond !A \in
  \Gamma$ 当且仅当存在某个完全 $\Sigma$ 一致集
  $\Delta$ 使得 $\Diamond\Delta \subseteq
  \Gamma$ 且 $!A \in \Delta$。
\end{prop}

\begin{proof}
假设 $\Gamma$ 是完全 $\Sigma$ 一致的。
根据 \Dual{} 及封闭性，$\Diamond!A \in
\Gamma$ 当且仅当 $\lnot\Box\lnot !A \in \Gamma$。
由于 $\Gamma$ 是完全 $\Sigma$ 一致的，根据
\olref[ccs]{prop:ccs-properties}\olref[ccs]{prop:ccs-lnot}，$\lnot\Box\lnot !A \in \Gamma$ 当且仅当 $\Box\lnot !A \notin \Gamma$。
由 \olref{prop:box}，$\Box\lnot !A \notin \Gamma$ 当且仅当存在某个完全
$\Sigma$ 一致集 $\Delta$ 满足 $\Box^{-1}\Gamma \subseteq \Delta$ 且 $\lnot!A \notin \Delta$。
现在考虑任意这样的~$\Delta$。
根据 \olref{lem:box-iff-diamond}，$\Box^{-1}\Gamma \subseteq \Delta$ 当且仅当
$\Diamond\Delta \subseteq \Gamma$。
同时，由
\olref[ccs]{prop:ccs-properties}\olref[ccs]{prop:ccs-lnot}，$\lnot !A \notin \Delta$ 当且仅当
$!A \in \Delta$。
因此，$\Diamond!A \in \Gamma$ 当且仅当存在某个完全 $\Sigma$ 一致集
$\Delta$ 使得 $\Diamond\Delta \subseteq \Gamma$ 且 $!A \in \Delta$。
\end{proof}

}{}}{}


\begin{prob}
  证明：若 $\Gamma$ 是完全 $\Sigma$ 一致的，则
  \iftag{prvBox}{$\Diamond !A \in \Gamma$ 当且仅当存在一个
    完全 $\Sigma$ 一致集 $\Delta$ 使得 $\Box^{-1}\Gamma
    \subseteq \Delta$ 且 $!A \in \Delta$。}{$\Box !A \in
    \Gamma$ 当且仅当对每个满足 $\Box^{-1} \Gamma \subseteq \Delta$ 的完全 $\Sigma$ 一致集
    $\Delta$，都有 $!A \in \Delta$。} \emph{请在不使用
    \olref[nml][com][mod]{lem:box-iff-diamond} 的前提下完成此题。}
\end{prob}

\end{document}
